\scanpage{19}
\uline{Définition de $\beta$}\qquad On définit
\[
\beta(\Pi,\widetilde\Pi,p,\mathcal G,\unsure{\sim|a})=(\unsure{\mathcal E|\Sigma},\mathcal G,\pi,\mathfrak Z)\quad\text{ainsi :}
\]
$\unsure{\mathcal E|\Sigma}=$ groupe des couples $(g,\bar g)$, avec $g\in\mathcal G$, $\bar g\in\operatorname{Aut}(\widetilde\Pi)$ compatible avec $g$
\[
\pi=\operatorname{Aut}(\widetilde\Pi/\Pi)
\]
$\unsure{\mathcal E|\Sigma}\longrightarrow\mathcal G$ par $(g,\bar g)\longmapsto g$, d'où
\diagram{\begin{tikzcd}[ampersand replacement=\&]
1 \arrow[r,dash] \& \pi \arrow[r] \& \unsure{\mathcal E|\Sigma} \arrow[r,dash] \& \mathcal G \arrow[r] \& 1
\end{tikzcd}}
$\mathfrak Z=\operatorname{Ob}\widetilde\Pi$, avec opération évidente de $\unsure{\mathcal E|\Sigma}$ sur $\mathfrak Z\ldots$

Variantes profinies\ldots

\textcircled{8} Considérons, avec notations \deleted{\unsure{du|de}} \S 7,
\[
\left(\diagram{\begin{tikzcd}[ampersand replacement=\&]
1 \arrow[r,dash] \& \pi \arrow[r,dash] \& \unsure{\mathcal E|\Sigma} \arrow[r] \& \mathcal G \arrow[r,dash] \& 1
\end{tikzcd}}
\quad\text{et }\mathfrak Z\quad\unsure{\mathcal E|\Sigma}\text{-ens}\ne\varnothing\text{ sur lequel }\pi\text{ opère librement}\right)
\]
$\longleftrightarrow(\Pi,\widetilde\Pi,p,\mathcal G,\unsure{\sim|a})$. On va interpréter la donnée d'un triple $(\inserted{\chi,}\unsure{\tau|\mathcal T},T,\Phi)$, \inserted{$\chi:$} où $\unsure{\tau|\mathcal T}$ est un $\unsure{\mathcal E|\Sigma}$-ens, $T$ un groupe $(\simeq\widehat{\mathbf Z}\;?)$, \inserted{$\chi:\mathcal G\longrightarrow\operatorname{Aut}(T)$ \unsure{$(\simeq\widehat{\mathbf Z}^{*}\unsure{)|?}\;\text{canon.}?)$}} et
\diagram{\begin{tikzcd}[ampersand replacement=\&]
\Phi:\unsure{\tau|\mathcal T} \arrow[r] \& \operatorname{Hom}_{\mathrm{gr}}(T,\pi)
\end{tikzcd}}
compatible avec op. de $\unsure{\mathcal E|\Sigma}$ (opérant
\mspage{18}
